\documentclass{article}
\usepackage{graphicx} % Required for inserting images
\usepackage{amssymb}
\usepackage{amsmath}

\newcommand{\bv}{\mathbf{v}}
\newcommand{\bu}{\mathbf{u}}
\newcommand{\bw}{\mathbf{w}}
\newcommand{\bO}{\mathbf{0}}
\newcommand{\R}{\mathbb{R}}
\newcommand{\Q}{\mathbb{Q}}
\newcommand{\C}{\mathbb{C}}
\newcommand{\F}{\mathbb{F}}

\title{Fundamentals Chapter 8}
\author{Siddharth K}
\date{November 2024}

\begin{document}

\maketitle

\textbf{1.} A nice basis for $V$ would be $\{1, x, \ldots, x^n\}$. The corresponding dual basis in $V^*$ would be $\{\delta_{1k}, \ldots, \delta_{nk}\}$ where each $\delta_{ik}$ takes in a polynomial of degree $k$ and returns $1$ if $k = i$ and $0$ otherwise. A nice basis in $V^*$ might look like $\{f_1, \ldots, f_n\}$ where $f_i(p(x)) = a_i$, the coefficient of $x^i$ in $p(x)$. \\

\textbf{3a.} Any element of $\bigcup \F[x]_{\le n}$ is a polynomial of finite degree, so it is in $\F[x]$. Also, any element of $\F[x]$ is a polynomial of degree $n$ for some $n$, so it is in $\bigcup \F[x]_{\le n}$. So the two sets are equal. \\

\textbf{3b.} Let $p(x) = \sum_{i = 0}^n a_ix^i$ and $q(x) = \sum_{i = 1}^n b_ix^i$ where $n$ is the maximum of the degrees of $p$ and $q$. We see that
$$\phi(p(x) + q(x)) = \phi\left(\sum_{i = 0}^n (a_i + b_i)x^i\right) = \sum_{i = 0}^n a_i + \sum_{i = 0}^n b_i = \phi(p(x)) + \phi(q(x)),$$
and
$$\phi(cp(x)) = \phi\left(\sum_{i = 0}^n ca_ix^i\right) = c\sum_{i = 0}^n a_i = c\phi(p(x)).$$

\textbf{3c.} Consider the function $\phi \in \F[x]^*$ where $\phi(p(x)) = p(1) = \sum_{i = 0}^\infty a_i$. For contradiction, assume that $\phi \in (\F[x]_{\le n})^*$ for some $n$. Then $\phi$ would only depend on the coefficients up to $a_n$, meaning $\phi(p(x)) = \sum_{i = 0}^n a_i$. But $\phi$ sums all the coefficients, and not just up to $n$. So no such finite $n$ will have the same return values as $\phi$ for all inputs, meaning this function cannot be equal to $\phi$. So $\phi \in \F[x]^*$ but not $\bigcup \F[x]_{\le n}$. \\

\textbf{4.} To show that $\varphi_\alpha \in V^*$ are linearly independent, we need to show that if $\sum_{q \in \Q} a_q\varphi_\alpha(\bv_q) = 0$ then $a_q = 0$ for all $q \in \Q$. We see that the sum is only $0$ if each $\varphi_\alpha(\bv_q)$ is always $0$, or if all $a_q = 0$. The former is clearly not always true, so $a_q = 0$ and $\varphi_\alpha$ are linearly independent. \\

\textbf{7.} Since a convex polygon is defined as the convex hull of finitely many points, by definition it can be expressed as a convex combination of its vertices. \\

\textbf{8.} The proof of the lemma left off by showing that if $\varphi(\bv_i) \le 1$ for all vertices $\bv_i$, then $$\varphi(s) = \varphi\left(\sum_{i = 0}^n \lambda_i \bv_i \right) \le 1$$ for all $s \in S$. This is what our goal was to prove, so $S^*$ is completely determined by the vertices of $S$. \\

\textbf{9.} Let's say the dual basis is $\{x, y\}$. In the $xy$-plane, the dual is the region satisfying the inequalities:
$$-\frac{1}{2}x-\frac{1}{2}y\le1,$$
$$-\frac{1}{2}x+\frac{3}{2}y\le1,$$
$$\frac{3}{2}x+\frac{3}{2}y\le1,$$
$$\frac{3}{2}x-\frac{1}{2}y\le1,$$
which forms a trapezoid with vertices $(0, -2), (-2, 0), (0, 2/3), (2/3, 0)$. \\

\textbf{10.} If the dual basis is $\{x, y\}$, the dual is the region in the $xy$-plane satisfying:
$$\pm x + y \le 1,$$
$$\pm(2x + y) \le 1,$$
which forms a triangle with vertices $(0, 1), (-2, 3), (-2/3, 1/3)$. \\

\textbf{11.} Consider the equilateral triangle with vertices $(2, 0), (-1, \sqrt{3}), (-1, -\sqrt{3})$. Its dual has vertices at $\left(\frac{1}{2}, \frac{\sqrt{3}}{2}\right), (-1, 0), \left(\frac{1}{2}, -\frac{\sqrt{3}}{2}\right)$ which lie on the original triangle. \\

\textbf{13.} The dual of the cube is the set of points that satisfy $\pm x \pm y \pm z \le 1$, which is an octahedron. The vertices of the octahedron meet the centers of the faces of the cube.

\end{document}
